\chapter{Social Choice}\label{ch:b2:social-choice}

Voting theory has a long list of things that can go wrong. With three or more candidates, the majority
may prefer $A$ to $B$, $B$ to $C$ and $C$ to $A$, so that no candidate beats every other; a candidate can lose
by gaining votes; a losing candidate's presence can change the winner. These are theorems about what
\emph{can} happen. The chart shows how often each happened in 182 American elections of 2004--2022 that used
ranked ballots, in which every voter lists the candidates in order.

\begin{openingchart}{Share of 182 US ranked-choice elections, 2004--2022, in which each flaw of the
  voting rule occurred, in per cent. All elections went to at least a second round of counting. The first
  bar is one election out of 182.}{Graham-Squire and McCune (2023), Table~7}
\begin{tikzpicture}
\begin{axis}[omrecord, xbar, bar width=9pt, xmin=0, xmax=5, ymin=0.4, ymax=7.6, y dir=reverse,
    ytick=data, yticklabels from table={figdata/bachelor-2/09-social-choice/out/flaws.csv}{flaw},
    xlabel={Per cent of elections}, grid=none, xmajorgrids, width=0.6\linewidth,
    height=0.48\linewidth, nodes near coords,
    every node near coord/.append style={font=\footnotesize}]
  \addplot[fill=omDef!70, draw=omDef] table[x=rate, y=pos]
    {figdata/bachelor-2/09-social-choice/out/flaws.csv};
\end{axis}
\end{tikzpicture}
\end{openingchart}

% ledger: b2:gsm2023-t7
% figdata: flaws_summary (no_winner 0.55, largest 3.8)
All but one of the 182 elections had a candidate who would have beaten every other one-to-one: 0.55 per
cent had none. The other flaws, each defined later in this chapter or in its exercises, occurred in at
most 3.8 per cent. The study finds one exception: in 52.5 per cent of the elections, the final-round
winner had a majority only of the ballots still counted, not of all ballots cast, because many voters
had not ranked the last two candidates.

\Cref{ch:b2:welfare-theorems} showed that markets reach efficient \omterm{def:b2:exchange-economies:allocation}{allocations}, and that choosing among
efficient \omterm{def:b2:exchange-economies:allocation}{allocations} needs a judgement the theorems do not supply. This chapter studies how such
judgements might be made: by comparing gains and losses, by a \omterm{def:b2:social-choice:swf}{social welfare function}, or by a vote. Each
method has a precise weakness, and the strongest result, Arrow's theorem, says that no voting rule escapes
all of them. The chart is the other half of the story: how often the weaknesses bite depends on what
voters' preferences actually look like.

\begin{recall}
From the school volumes: majority voting and the idea that collective decisions are made by rules. From
\cref{ch:b2:exchange-economies,ch:b2:welfare-theorems}: Pareto efficiency, the \omterm{def:b2:exchange-economies:edgeworth-box}{Edgeworth box}, and the
second welfare theorem's \omterm{def:b2:welfare-theorems:lump-sum}{lump-sum transfers}.
\end{recall}

\begin{notation}[Symbols of this chapter]\label{not:b2:social-choice:symbols}
Voters or citizens are $i = 1, \dots, n$; alternatives are $x, y, z$ (candidates, budgets, policies), and
$m$ is their number. $x \succ_i y$ means voter $i$ ranks $x$ above $y$; a profile is a list of every voter's
ranking. $\succsim$ is a social ordering and $\succ$ its strict part. $u_i$ is person $i$'s utility and
$W(u_1, \dots, u_n)$ a \omterm{def:b2:social-choice:swf}{social welfare function}.
\end{notation}

\section{The Pareto criterion and compensation tests}\label{sec:b2:social-choice:pareto}

\begin{definition}[Pareto improvement]\label{def:b2:social-choice:pareto-improvement}
A change from one state of the economy to another is a \emph{Pareto improvement}\index{Pareto
improvement} if it makes at least one person better off and no one worse off.
\end{definition}

The Pareto criterion is the least controversial way of ranking states: almost anyone accepts that a
\omterm{def:b2:social-choice:pareto-improvement}{Pareto improvement} is an improvement. Its weakness is that it is silent about most comparisons. Almost
every policy, a new road, a tariff cut, a tax, makes someone worse off, and the criterion then has nothing
to say. A \omterm{def:b2:social-choice:compensation-test}{compensation test} widens its reach.

\begin{definition}[Compensation test]\label{def:b2:social-choice:compensation-test}
A change from state $A$ to state $B$ passes the \emph{compensation test}\index{compensation test}
(Kaldor--Hicks test) if the people who gain from it could compensate the people who lose and still be
better off: some reallocation of what $B$ makes possible is a \omterm{def:b2:social-choice:pareto-improvement}{Pareto improvement} on $A$. The compensation
need not be paid.
\end{definition}

The test is what lies behind adding up gains and losses in money, as in cost--benefit analysis: if gains
exceed losses, the winners could compensate the losers. It has two weaknesses. Because compensation is
not paid, the losers stay losers, so passing the test is not a \omterm{def:b2:social-choice:pareto-improvement}{Pareto improvement}. And the test can approve
a change and its reverse.

Two policies, $A$ and $B$, lead to different sets of possible utilities for two people. Under $A$ the
utility possibility frontier is $u_1 + 2u_2 = 10$; under $B$ it is $2u_1 + u_2 = 10$. State $A$ gives
$(u_1, u_2) = (2, 4)$; state $B$ gives $(4, 2)$. From $A$, moving to $B$ passes the test, because policy $B$ allows
$(3, 4)$, which is better for both than $(2, 4)$. From $B$, moving back to $A$ also passes, because policy $A$
allows $(4, 3)$, which is better for both than $(4, 2)$. The test approves both moves.

\begin{omfigure}
\begin{tikzpicture}[econ]
\begin{axis}[omrecord, grid=none, xmin=0, xmax=6.5, ymin=0, ymax=6.5, xlabel={$u_1$}, ylabel={$u_2$},
    xtick={0,2,4,6}, ytick={0,2,4,6}, width=0.5\linewidth, height=0.5\linewidth]
  \addplot[omDef, very thick, domain=0:6.5, samples=2] {(10 - x)/2};
  \addplot[omThm, very thick, domain=1.75:5, samples=2] {10 - 2*x};
  \node[eqpoint] at (axis cs:2,4) {};
  \node[anchor=north east, font=\footnotesize] at (axis cs:1.95,3.95) {$A$};
  \node[eqpoint] at (axis cs:4,2) {};
  \node[anchor=north, font=\footnotesize] at (axis cs:3.9,1.9) {$B$};
  \node[draw, circle, inner sep=1.5pt] at (axis cs:3,4) {};
  \node[anchor=south west, font=\footnotesize] at (axis cs:3,4) {$(3, 4)$};
  \node[draw, circle, inner sep=1.5pt] at (axis cs:4,3) {};
  \node[anchor=south west, font=\footnotesize] at (axis cs:4,3) {$(4, 3)$};
  \node[anchor=west, font=\footnotesize, omDef] at (axis cs:4.6,2.4) {policy $A$};
  \node[anchor=west, font=\footnotesize, omThm] at (axis cs:2.4,5.6) {policy $B$};
\end{axis}
\end{tikzpicture}

{\small A compensation-test reversal. The two utility possibility frontiers cross. From $A$, policy $B$'s
frontier passes above and to the right of $A$, at $(3, 4)$; from $B$, policy $A$'s frontier passes above
and to the right of $B$, at $(4, 3)$. Each state passes the test against the other.}
\end{omfigure}

The reversal needs the frontiers to cross, which they can when the two policies change what is produced
and so whose tastes it favours. When they do not cross, the test gives a consistent answer, and it is then
equivalent to asking whether one frontier lies outside the other: a statement about efficiency, not about
the distribution actually reached.

\section{Social welfare functions as tools}\label{sec:b2:social-choice:swf}

\begin{definition}[Social welfare function]\label{def:b2:social-choice:swf}
A \emph{social welfare function}\index{social welfare function} $W(u_1, \dots, u_n)$ ranks states of the
economy by the utilities they give to its $n$ members; it increases in each $u_i$. Common forms are the
utilitarian $W = \sum_i u_i$, the maximin $W = \min_i u_i$, and the weighted sum $W = \sum_i a_i u_i$ with
weights $a_i > 0$.
\end{definition}

A \omterm{def:b2:social-choice:swf}{social welfare function} does not come from economics; it states a judgement, which economics can then
apply consistently. Every one that increases in each $u_i$ respects the Pareto criterion, and each makes a
different trade-off when the criterion is silent. A worked case shows how far apart they can be.

% code: exercises/solutions.py transfers
Two people have incomes of 80 and 20. Each dollar taken from the first reaches the second as 80 cents: the
rest is lost in administration and in changed behaviour. With utilities equal to income, the utilitarian
function forbids any transfer, since each one lowers the sum. With utilities $\ln m$, the utilitarian
optimum sets $1/(80 - t) = 0.8/(20 + 0.8t)$, so $t = 27.5$: incomes become 52.5 and 42. The maximin
function equalises incomes, $80 - t = 20 + 0.8t$, so $t = 33.33$ and both have 46.67, although 6.67 dollars
are lost on the way. The three answers come from the same facts and different judgements. What economics
can supply is the facts: how large the leak is, which depends on how people respond to taxes and
transfers, a question for the public economics of the Year 3 volume.

\begin{definition}[Envy-free allocation]\label{def:b2:social-choice:envy-free}
An \omterm{def:b2:exchange-economies:allocation}{allocation} is \emph{envy-free}\index{envy-free allocation} if every person weakly prefers her own bundle
to the bundle of every other person.
\end{definition}

\begin{remark}
Envy-freeness compares bundles rather than utilities, so it needs no comparison of one person's
satisfaction with another's. In an \omterm{def:b2:exchange-economies:exchange-economy}{exchange economy}, the \omterm{def:b2:exchange-economies:competitive-equilibrium}{competitive equilibrium} from equal \omterm{def:b2:exchange-economies:endowment}{endowments} is
both \omterm{def:b2:social-choice:envy-free}{envy-free} and \omterm{def:b2:exchange-economies:pareto-efficient}{Pareto-efficient}: everyone faces the same budget, so anyone who preferred another's
bundle could have bought it.
\end{remark}

\section{Voting cycles}\label{sec:b2:social-choice:cycles}

When a group decides by voting, it does not need to compare utilities: each voter ranks the
alternatives, and a rule turns the rankings into a decision. The simplest rule for two alternatives,
majority voting, has every good property one could ask for. With three, trouble starts.

\begin{definition}[Condorcet winner]\label{def:b2:social-choice:condorcet-winner}
A \emph{Condorcet winner}\index{Condorcet winner} is an alternative that a majority of voters prefers to
each other alternative, in one-to-one comparisons.
\end{definition}

Three voters rank three alternatives: voter 1 has $x \succ y \succ z$, voter 2 has $y \succ z \succ x$,
voter 3 has $z \succ x \succ y$. Voters 1 and 3 prefer $x$ to $y$; voters 1 and 2 prefer $y$ to $z$;
voters 2 and 3 prefer $z$ to $x$. Each comparison is won two votes to one, and the majority's preferences go
round in a circle.

\begin{definition}[Condorcet paradox]\label{def:b2:social-choice:condorcet-paradox}
The \emph{Condorcet paradox}\index{Condorcet paradox} (voting cycle) is the possibility that majority
voting over pairs is intransitive: a majority prefers $x$ to $y$, a majority $y$ to $z$ and a majority $z$
to $x$, although every voter's ranking is transitive. There is then no \omterm{def:b2:social-choice:condorcet-winner}{Condorcet winner}.
\end{definition}

\begin{omfigure}
\begin{tikzpicture}[econ, font=\footnotesize, alt/.style={draw, circle, minimum size=8mm,
    fill=omDef!15}]
  \node[alt] (x) at (90:1.6) {$x$};
  \node[alt] (y) at (210:1.6) {$y$};
  \node[alt] (z) at (330:1.6) {$z$};
  \draw[->, very thick, omThm] (x) -- node[left=2pt] {2 to 1} (y);
  \draw[->, very thick, omThm] (y) -- node[below=2pt] {2 to 1} (z);
  \draw[->, very thick, omThm] (z) -- node[right=2pt] {2 to 1} (x);
\end{tikzpicture}

{\small A majority cycle. An arrow from one alternative to another means that a majority prefers the
first. Whatever is chosen, a majority prefers something else.}
\end{omfigure}

With a cycle, the outcome depends on the order of the votes. If $x$ first meets $y$ and the winner then
meets $z$, $z$ wins; with the agenda $y$ against $z$, then $x$, $x$ wins. Whoever sets the agenda chooses the
outcome (\cref{pb:b2:social-choice:1}, Part IV).

How likely is a cycle? If every voter's ranking were drawn at random, with all orders equally likely, the
chance of no \omterm{def:b2:social-choice:condorcet-winner}{Condorcet winner} in a large electorate can be computed. The listing estimates it by
simulation, for 501 voters and three to five candidates, and repeats the count for \omterm{def:b2:social-choice:single-peaked}{single-peaked
preferences} (\cref{sec:b2:social-choice:median}), where each voter has a favourite point on a line and
likes candidates less the further they are from it.

\omcode{code/bachelor-2/09-social-choice/listings/condorcet_random.py}{1}{27}{Counting elections without
a \omterm{def:b2:social-choice:condorcet-winner}{Condorcet winner}, for random and for \omterm{def:b2:social-choice:single-peaked}{single-peaked preferences}.}\label{lst:b2:social-choice:random}

\omoutput{code/bachelor-2/09-social-choice/outputs/condorcet_random.txt}{With random preferences cycles are
common, and they grow with the number of candidates; with \omterm{def:b2:social-choice:single-peaked}{single-peaked preferences} there are
none.}\label{out:b2:social-choice:random}

% code: listings/condorcet_random.py
With random rankings, 9.2 per cent of three-candidate elections and 24.2 per cent of five-candidate
elections have no \omterm{def:b2:social-choice:condorcet-winner}{Condorcet winner}. In the ranked-choice elections of the opening chart, which often had
more than three candidates, the share was 0.55 per cent. Real preferences are far from random: they share
structure, and the structure is what makes majority voting work.

\section{Arrow's theorem}\label{sec:b2:social-choice:arrow}

Majority voting over pairs can cycle. Is there a better rule? A rule should turn any profile of rankings
into a ranking for society.

\begin{definition}[Social preference ordering]\label{def:b2:social-choice:social-ordering}
A \emph{social preference ordering}\index{social preference ordering} is a complete and transitive
ranking of the alternatives on behalf of society. A rule that assigns one to each profile of the voters'
rankings satisfies \emph{unrestricted domain} if it is defined for every possible profile, and the
\emph{weak Pareto} condition if society ranks $x$ above $y$ whenever every voter does.
\end{definition}

\begin{definition}[Independence of irrelevant alternatives]\label{def:b2:social-choice:iia}
A rule satisfies \emph{independence of irrelevant alternatives}\index{independence of irrelevant
alternatives} if society's ranking of any two alternatives $x$ and $y$ depends only on how each voter ranks
$x$ against $y$, not on where voters place other alternatives.
\end{definition}

\begin{definition}[Dictator]\label{def:b2:social-choice:dictator}
A voter is a \emph{dictator}\index{dictator (social choice)} for a rule if, for every profile and every pair
$x, y$, society ranks $x$ above $y$ whenever she does.
\end{definition}

\begin{theorem}[Arrow's impossibility theorem]\label{thm:b2:social-choice:arrow}
With at least three alternatives and voters with strict rankings, every rule that produces a \omterm{def:b2:social-choice:social-ordering}{social
preference ordering} and satisfies unrestricted domain, weak Pareto and \omterm{def:b2:social-choice:iia}{independence of irrelevant
alternatives} has a \omterm{def:b2:social-choice:dictator}{dictator}.
\end{theorem}

\begin{proof}
\emph{Step 1: an alternative that every voter puts at an extreme is at an extreme socially.} Let every voter
place $b$ at the top or the bottom of her ranking, and suppose society does not: there are $a$ and $c$ with
$a \succsim b$ and $b \succsim c$. Change every voter's ranking by moving $c$ just above $a$. This changes no
voter's comparison of $a$ with $b$ or of $c$ with $b$, since $b$ is at an extreme, so by independence $a
\succsim b$ and $b \succsim c$ still hold, and by transitivity $a \succsim c$. But every voter now prefers $c$
to $a$, so by weak Pareto $c \succ a$: a contradiction.

\emph{Step 2: a pivotal voter.} Start from a profile where every voter places $b$ at the bottom; by weak
Pareto, society does too. Move $b$ to the top of the voters' rankings one at a time, voter 1 first. When all
have moved it, society places $b$ at the top. By step 1, at every stage $b$ is at the top or the bottom
socially, so there is a first voter $n^*$ whose move takes it from the bottom to the top. Call the profiles
just before and just after her move I and II.

\emph{Step 3: $n^*$ decides every pair without $b$.} Take $a, c \ne b$. Starting from II, let $n^*$ move $a$
above $b$, so that she ranks $a \succ b \succ c$, and let the other voters rearrange $a$ and $c$ in any way,
leaving $b$ where it was. Call this profile III. In III every voter compares $a$ and $b$ as in I, where $b$
was at the bottom socially, so $a \succ b$. Every voter compares $b$ and $c$ as in II, where $b$ was at the
top, so $b \succ c$. Hence $a \succ c$ in III. By independence, society ranks $a$ above $c$ whenever $n^*$
does, whatever the others do.

\emph{Step 4: $n^*$ also decides the pairs with $b$.} Take $c \ne b$. Steps 2 and 3, run with $c$ in the role of
$b$, give a voter who decides every pair without $c$, among them every pair $(a, b)$. Between profiles I and
II only $n^*$'s ranking changed, and society's ranking of $b$ against $a$ changed with it; so that voter
must be $n^*$. Hence $n^*$ is a \omterm{def:b2:social-choice:dictator}{dictator}.
\end{proof}

The theorem does not say that every collective decision is arbitrary; it says that no rule keeps all four
conditions on every possible profile. Something has to give. The ranked-choice elections of the opening
chart use instant runoff, which breaks independence; majority voting over pairs breaks the requirement of
a transitive ordering, but only for profiles that cycle. The record suggests where the escape lies:
real profiles rarely come from the whole unrestricted domain.

\section{Single-peaked preferences and the median voter}\label{sec:b2:social-choice:median}

\begin{definition}[Single-peaked preferences]\label{def:b2:social-choice:single-peaked}
Alternatives are ordered on a line, such as budgets from low to high. A voter's preferences are
\emph{single-peaked}\index{single-peaked preferences} if she has a favourite alternative and, on each side
of it, likes alternatives less the further they are from it.
\end{definition}

\begin{omfigure}
\begin{tikzpicture}[econ]
\begin{axis}[omrecord, grid=none, xmin=0.5, xmax=3.5, ymin=0, ymax=3.5, xtick={1,2,3},
    xticklabels={$L$, $M$, $R$}, ytick=\empty, xlabel={Alternatives on a line},
    ylabel={Ranking (higher is better)}, width=0.5\linewidth, height=0.42\linewidth,
    legend style={at={(1.04,0.5)}, anchor=west, font=\footnotesize}]
  \addplot[omDef, very thick, mark=*] coordinates {(1,3) (2,2) (3,1)};
  \addplot[black!60, thick, mark=square*] coordinates {(1,1) (2,3) (3,2)};
  \addplot[omThm, thick, dashed, mark=triangle*] coordinates {(1,2) (2,1) (3,3)};
  \legend{$L \succ M \succ R$, $M \succ R \succ L$, $R \succ L \succ M$}
\end{axis}
\end{tikzpicture}

{\small The voters of the cycle in \cref{sec:b2:social-choice:cycles}, with $x = L$, $y = M$, $z = R$. The first
two rankings are \omterm{def:b2:social-choice:single-peaked}{single-peaked}. The third has two peaks, at $L$ and $R$, with a dip at $M$: it is not.}
\end{omfigure}

\begin{definition}[Median voter theorem]\label{def:b2:social-choice:median-voter}
The \emph{median voter theorem}\index{median voter theorem} is the result that, with \omterm{def:b2:social-choice:single-peaked}{single-peaked
preferences} and an odd number of voters, the favourite of the median voter wins every majority vote over
pairs (\cref{thm:b2:social-choice:median-voter}).
\end{definition}

\begin{theorem}[Median voter theorem]\label{thm:b2:social-choice:median-voter}
Let the number of voters $n$ be odd and every voter's preferences be \omterm{def:b2:social-choice:single-peaked}{single-peaked} on the same line. Let $x_m$
be the median of the voters' favourite alternatives. Then $x_m$ is a \omterm{def:b2:social-choice:condorcet-winner}{Condorcet winner}: a majority prefers it
to every other alternative.
\end{theorem}

\begin{proof}
Take any $y < x_m$. At least $(n + 1)/2$ voters have a favourite at or above $x_m$. Each of them is on the
side of $x_m$ away from $y$, so by single-peakedness she prefers $x_m$, which is closer to her peak along the
line, to $y$. That is a majority. The case $y > x_m$ is symmetric, using the $(n + 1)/2$ voters with favourites
at or below $x_m$.
\end{proof}

Single-peakedness rules out the cycle: in the proof, any alternative loses to the median's favourite, so
no chain of majorities can go round. It also says which alternative wins: neither the average voter's
favourite nor the one that maximises a utilitarian sum, but the median's. If those who favour high budgets
care much more than those who favour low ones, majority voting ignores the intensity.

\begin{model}{The median-voter model}
\assumes Alternatives on one dimension; \omterm{def:b2:social-choice:single-peaked}{single-peaked preferences}; majority voting over pairs; every
voter votes.
\predicts The median voter's favourite wins; changes in the electorate move the outcome only through the
median; candidates who compete for votes converge on the median.
\tested Equations for local public spending built on the median voter's income and tax share fit municipal
data, which is consistent with the model but does not single it out: a test needs cases where the median
voter's favourite differs from what an agenda setter or the mean voter would choose. The Year 3 volume
reviews that evidence.
\breaks Several dimensions at once (a budget and its use); turnout that differs across voters; whoever
sets the agenda; intensity of preferences.
\end{model}

\section{Strategic voting}\label{sec:b2:social-choice:strategic}

All the rules so far took the rankings as given. A voter who knows the rule may report a ranking that is not
her own, if that brings a better result.

\begin{definition}[Strategy-proof voting rule]\label{def:b2:social-choice:strategy-proof}
A voting rule is \emph{strategy-proof}\index{strategy-proof voting rule} if no voter can ever obtain an
outcome she prefers by reporting a ranking other than her true one, whatever the others report.
\end{definition}

% code: exercises/solutions.py borda
The Borda count gives each alternative $m - 1$ points for a first place, $m - 2$ for a second, and so on.
Three voters rank $x \succ y \succ z$ and two rank $y \succ z \succ x$. Then $x$ is the \omterm{def:b2:social-choice:condorcet-winner}{Condorcet winner}, since three
of the five prefer it to each other alternative, but the Borda count gives $x$ six points, $y$ seven and $z$ two:
$y$ wins. If the three supporters of $x$ report $x \succ z \succ y$ instead, $y$ falls to four points and $z$ rises to
five, and $x$ wins with six. By moving $y$ to the bottom, they changed the result in their favour.

The Gibbard--Satterthwaite theorem says that this is no accident of the Borda count. With at least three
alternatives, every voting rule that is \omterm{def:b2:social-choice:strategy-proof}{strategy-proof} and can produce each alternative as the winner
has a \omterm{def:b2:social-choice:dictator}{dictator}. \omterm{def:b2:social-choice:single-peaked}{Single-peaked preferences} again offer the way out: with them, choosing the median of
the reported favourites is \omterm{def:b2:social-choice:strategy-proof}{strategy-proof}, since a voter who reports a favourite on the other side of the
median moves the outcome away from her (\cref{pb:b2:social-choice:1}, question 16).

\section{Exercises}\label{sec:b2:social-choice:exercises}

\begin{exercise}[$\star$]\label{exo:b2:social-choice:1}
Five voters rank three alternatives: two rank $A \succ B \succ C$, two rank $B \succ C \succ A$, and one ranks
$C \succ A \succ B$. Find the result of each one-to-one majority vote. Is there a \omterm{def:b2:social-choice:condorcet-winner}{Condorcet winner}?
\end{exercise}

\begin{exercise}[$\star$]\label{exo:b2:social-choice:2}
Three states give two people the utilities $x = (5, 3)$, $y = (6, 3)$ and $z = (8, 1)$. Which comparisons can
the Pareto criterion decide?
\end{exercise}

\begin{exercise}[$\star$]\label{exo:b2:social-choice:3}
Read the opening chart. In how many of the 182 elections was there no \omterm{def:b2:social-choice:condorcet-winner}{Condorcet winner}? Compare with the
listing's estimate for random preferences, and say what the difference suggests.
\end{exercise}

\begin{exercise}[$\star$]\label{exo:b2:social-choice:4}
A town decides its library budget by majority vote over pairs. Describe a result that would refute the
\omterm{def:b2:social-choice:median-voter}{median voter theorem}'s prediction, and what you would need to check before concluding that it did.
\end{exercise}

\begin{exercise}[$\star\star$]\label{exo:b2:social-choice:5}
Three voters rank $A \succ B \succ C$ and two rank $B \succ C \succ A$. Find the \omterm{def:b2:social-choice:condorcet-winner}{Condorcet winner} and the Borda
winner. Which property does the Borda count violate here?
\end{exercise}

\begin{exercise}[$\star\star$]\label{exo:b2:social-choice:6}
On the line $L < M < R$, which of the rankings $L \succ M \succ R$, $M \succ R \succ L$ and $R \succ L \succ M$ are
\omterm{def:b2:social-choice:single-peaked}{single-peaked}? Explain why a \omterm{def:b2:social-choice:single-peaked}{single-peaked} voter never ranks the middle alternative last.
\end{exercise}

\begin{exercise}[$\star\star$]\label{exo:b2:social-choice:7}
Five voters have favourite spending levels 10, 20, 30, 40 and 50. Two voters with favourites 60 and 70 join
the electorate. Predict the direction of the change in the outcome, then compute it. Would two new voters
with favourites 5 and 70 change it?
\end{exercise}

\begin{exercise}[$\star\star$]\label{exo:b2:social-choice:8}
Policy $A$ allows the utilities $u_1 + 3u_2 = 12$, policy $B$ allows $3u_1 + u_2 = 12$. State $A$ is $(1.5, 3.5)$
and state $B$ is $(3.5, 1.5)$. Show that the \omterm{def:b2:social-choice:compensation-test}{compensation test} approves the move from $A$ to $B$ and the move from $B$
to $A$.
\end{exercise}

\begin{exercise}[$\star\star$]\label{exo:b2:social-choice:9}
Towns vote on local spending in referenda. Design a test of the median voter model with data on the
outcomes and on the residents' incomes, and say what result would count against it.
\end{exercise}

\begin{exercise}[$\star\star\star$]\label{exo:b2:social-choice:10}
In the proof of \cref{thm:b2:social-choice:arrow}: (a) explain why, in step 2, $b$ cannot stop halfway
up society's ranking; (b) in step 3, check that profile III can be built for any rearrangement of $a$ and
$c$ by the other voters; (c) why does step 4 need profiles I and II?
\end{exercise}

\begin{exercise}[$\star\star\star$]\label{exo:b2:social-choice:11}
Prove the \omterm{def:b2:social-choice:median-voter}{median voter theorem} for an even number of voters: every alternative between the two middle
favourites, inclusive, is beaten by no other alternative in a majority vote, although it may tie.
\end{exercise}

\begin{exercise}[$\star\star\star$]\label{exo:b2:social-choice:12}
In \cref{exo:b2:social-choice:5}, show that the three voters who rank $A$ first can make $A$ win the Borda count
by misreporting. Could the two others then answer with a misreport of their own? What does this say
about the rule?
\end{exercise}

\section{Problem: Brenmoor's park budget}\label{sec:b2:social-choice:problem}

\begin{problem}[{Weekend problem --- the seven councillors of Brenmoor, a city of Brenland, an invented
economy, choose a budget for a new park: which budget wins each kind of vote, and how does adding a second
question change the result}]
\label{pb:b2:social-choice:1}
The options are budgets of 200, 400, 600 and 800 thousand dalens. Each councillor has a favourite budget,
not necessarily one of the options, and ranks the options by their distance from it: 250, 380, 450, 520,
640, 690 and 900.

\medskip
\noindent\textbf{Part I --- Ballots.}
\begin{enumerate}
  \item Write each councillor's ranking of the four options.
  \item Which option wins a plurality vote (most first places)?
  \item Compute the Borda scores (3, 2, 1 and 0 points).
  \item Find the result of each one-to-one majority vote involving 600.
  \item Is there a \omterm{def:b2:social-choice:condorcet-winner}{Condorcet winner}?
  \item Do the three rules agree?
\end{enumerate}

\medskip
\noindent\textbf{Part II --- Single-peakedness.}
\begin{enumerate}[resume]
  \item Show that every ranking is \omterm{def:b2:social-choice:single-peaked}{single-peaked} on the line $200 < 400 < 600 < 800$.
  \item A councillor ranks $200 \succ 800 \succ 400 \succ 600$. Is this \omterm{def:b2:social-choice:single-peaked}{single-peaked}? What kind of view might
        produce it?
  \item Why does ranking options by distance from a favourite always give \omterm{def:b2:social-choice:single-peaked}{single-peaked preferences}?
  \item Which theorem now applies, and what does it predict?
\end{enumerate}

\medskip
\noindent\textbf{Part III --- The median and pairwise votes.}
\begin{enumerate}[resume]
  \item Find the median favourite budget.
  \item Which option does the median councillor rank first?
  \item Check that this option beats every other in one-to-one votes.
  \item The councillor whose favourite is 250 is replaced by one whose favourite is 950. Find the new median
        favourite and the new \omterm{def:b2:social-choice:condorcet-winner}{Condorcet winner}.
  \item Councillors 380 and 450 both report $400 \succ 200 \succ 800 \succ 600$ in a Borda count. Find the new
        scores. Could either of them change the result alone?
  \item Under pairwise majority voting, could any councillor obtain a better result by misreporting?
\end{enumerate}

\medskip
\noindent\textbf{Part IV --- A second question.} The council must also decide where the park goes.
Three proposals combine a budget and a site: $X$, $Y$ and $Z$. Three councillors rank $X \succ Y \succ Z$, two
rank $Y \succ Z \succ X$, and two rank $Z \succ X \succ Y$.
\begin{enumerate}[resume]
  \item Find the result of each one-to-one vote.
  \item Show that there is a cycle.
  \item The mayor first puts $X$ against $Y$, then the winner against $Z$. Which proposal is adopted?
  \item Find the outcome of each of the other two agendas.
  \item Which condition of Arrow's theorem do the rankings violate, and why can two questions produce
        this when one could not?
  \item Under the agenda of question 19, would the three councillors who rank $X$ first gain by voting
        for $Y$ in the first round?
  \item What does this suggest about who holds power in a council that votes over pairs?
  \item State the result: the budget that wins every pairwise vote, and the outcome of each agenda once the
        site is added.
\end{enumerate}
\end{problem}
